% Компиляция PDF: ./build_pdf.sh (см. docs/slides/README.md) или дважды: xelatex presentation.tex
\documentclass[aspectratio=169]{beamer}

\usepackage{polyglossia}
\setmainlanguage{russian}
\usepackage{fontspec}
\setmainfont{DejaVu Sans}
\setsansfont{DejaVu Sans}
\setmonofont{DejaVu Sans Mono}

% Kaspersky-like beamer styling (not official brand)
% Compile with lualatex or xelatex

\usepackage{fontspec}
\usepackage{xcolor}
\usepackage{graphicx}
\setbeamertemplate{navigation symbols}{}
% Язык задаётся в основном файле презентации (polyglossia: \setmainlanguage{russian})

% --- Fonts (choose installed fonts with Cyrillic support) ---
% Хорошие нейтральные варианты: Inter, IBM Plex Sans, Noto Sans, PT Sans.
% \setmainfont{Inter}[
% \setmainfont{TeX Gyre Termes}[
%   Extension      = .ttf,
% %   UprightFont    = *-Regular,
% %   ItalicFont     = *-Italic,
%   BoldFont       = *-Bold,
%   BoldItalicFont = *-BoldItalic
% ]
% \setsansfont{TeX Gyre Termes}[
%   Extension      = .ttf,
% %   UprightFont    = *-Regular,
%   ItalicFont     = *-Italic,
%   BoldFont       = *-Bold,
%   BoldItalicFont = *-BoldItalic
% ]
% \setmonofont{DejaVu Sans Mono}

% Если Inter не установлен, замените на Noto Sans:
% setmainfont{Noto Sans}
% setsansfont{Noto Sans}

% --- Color palette (approximate) ---
\definecolor{KaspGreen}{HTML}{00A88F} % акцент
\definecolor{KaspDark}{HTML}{0B0F14}  % почти чёрный
\definecolor{KaspGray}{HTML}{6B7280}  % серый текст
\definecolor{KaspLight}{HTML}{F5F7FA} % светлый фон

% --- Beamer base ---
\mode<presentation>
{
  \setbeamertemplate{navigation symbols}{}
  \setbeamersize{text margin left=8mm, text margin right=8mm}
}

% --- Global colors ---
\setbeamercolor{normal text}{fg=KaspDark,bg=white}
\setbeamercolor{frametitle}{fg=KaspDark,bg=white}
\setbeamercolor{title}{fg=KaspDark}
\setbeamercolor{subtitle}{fg=KaspGray}
\setbeamercolor{structure}{fg=KaspGreen}

% --- Typography ---
\setbeamerfont{title}{series=\bfseries,size=\LARGE}
\setbeamerfont{frametitle}{series=\bfseries,size=\Large}
\setbeamerfont{subtitle}{size=\large}
\setbeamerfont{footline}{size=\scriptsize}

% % --- Blocks ---
\setbeamercolor{block title}{fg=white,bg=KaspDark}
\setbeamercolor{block body}{fg=KaspDark,bg=KaspLight}
\setbeamercolor{block title alerted}{fg=white,bg=KaspGreen}
\setbeamercolor{block body alerted}{fg=KaspDark,bg=KaspLight}

% % --- Itemize ---
\setbeamertemplate{itemize item}{\color{KaspGreen}\large\textbullet}
\setbeamertemplate{itemize subitem}{\color{KaspGreen}\small\textbullet}

% --- Title page template ---
\setbeamertemplate{title page}{
  \vspace*{6mm}
  {\usebeamerfont{title}\inserttitle\par}
  \vspace*{8mm}
  {\color{KaspGray}\insertauthor\par}
  {\color{KaspGray}\insertinstitute\par}
  \vspace*{2mm}
  {\color{KaspGray}\insertdate\par}
  \vfill
  {\color{KaspGreen}\rule{\linewidth}{1.2pt}}
}

% --- Frame title with accent line ---
\setbeamertemplate{frametitle}{
  \vspace*{1mm}
  {\usebeamerfont{frametitle}\insertframetitle\par}
  \vspace*{1mm}
  {\color{KaspGreen}\rule{0.22\linewidth}{0.8pt}}
  \vspace*{1mm}
}

%--- Footline (left: section, right: slide number) ---
% \setbeamertemplate{footline}{
%   \hbox{
%     \begin{beamercolorbox}[wd=.78\paperwidth,ht=3.2ex,dp=1.4ex,left]{footline}
%       {\hspace*{8mm}\color{KaspGray}\insertsection}
%     \end{beamercolorbox}
%     \begin{beamercolorbox}[wd=.22\paperwidth,ht=3.2ex,dp=1.4ex,right]{footline}
%       {\color{KaspGray}\insertframenumber{} / \inserttotalframenumber\hspace*{8mm}}
%     \end{beamercolorbox}
%   }
% }

% В преамбуле после usetheme{...} (если есть)

% Убираем стандартные значки навигации
% \setbeamertemplate{navigation symbols}{}

% Нижний колонтитул
\setbeamertemplate{footline}{
  \leavevmode%
  \hbox{%
    \begin{beamercolorbox}[wd=.33\paperwidth,ht=2.5ex,dp=1.2ex,left]{author in head/foot}%
      \hspace*{1ex}\insertshortauthor
    \end{beamercolorbox}%
    \begin{beamercolorbox}[wd=.34\paperwidth,ht=2.5ex,dp=1.2ex,center]{title in head/foot}%
      \insertshorttitle
    \end{beamercolorbox}%
    \begin{beamercolorbox}[wd=.33\paperwidth,ht=2.5ex,dp=1.2ex,right]{date in head/foot}%
      \insertframenumber{} / \inserttotalframenumber\hspace*{1ex}
    \end{beamercolorbox}%
  }%
  \vskip0pt%
}


\title{Инструктаж перед стартом конкурса}
\subtitle{Что прочитать, что сделать, как не потерять баллы}
\author{Учебный стенд}
\institute{Конкурс / прототип}
\date{\today}

\begin{document}

\begin{frame}
  \titlepage
\end{frame}

\begin{frame}{Формат этого инструктажа}
  \begin{itemize}
    \item Слайды рассчитаны на \textbf{10--15 минут} до начала соревнований.
    \item Цель: дать минимально достаточный маршрут, чтобы сразу начать работу без лишнего чтения.
    \item Фокус: \textbf{3 обязательных документа}, практические шаги на 2 дня и контрольные проверки перед сдачей.
  \end{itemize}
\end{frame}

\begin{frame}{Что нужно прочитать в первую очередь}
  \begin{enumerate}
    \item \texttt{docs/contest\_task.md} --- постановка задачи и обязательные артефакты.
    \item \texttt{docs/quickstart\_2days.md} --- пошаговый маршрут на 2 дня.
    \item \texttt{docs/criteria\_rubric.md} --- как начисляются баллы по C01--C25.
  \end{enumerate}
  \vspace{0.6em}
  \textbf{Важно:} этого набора достаточно, чтобы начать выполнять задание.
\end{frame}

\begin{frame}{Вспомогательные документы (по необходимости)}
  \begin{itemize}
    \item \texttt{docs/README.md} --- карта документации.
    \item \texttt{docs/contest\_regulations.md} --- полный регламент и правила разрешения равенства баллов.
    \item \texttt{docs/certification\_process.md} --- подробная процедура сертификации.
    \item \texttt{docs/architecture.md}, \texttt{docs/context.md} --- архитектура и предметная область.
    \item \texttt{docs/sbom\_guide.md}, \texttt{docs/tara\_abu.md}, \texttt{docs/security\_tests.md}, \texttt{docs/faq\_participant.md}.
  \end{itemize}
  \vspace{0.4em}
  Рекомендованная последовательность уже дана в \texttt{quickstart\_2days.md}.
\end{frame}

\begin{frame}{О чём само задание (кратко)}
  \begin{itemize}
    \item Вы дорабатываете решение в \texttt{src\_solution/} на базе \texttt{src\_starting\_point/}.
    \item Ключевая идея: минимальная и проверяемая ДВБ, явные IPC-политики, контролируемое междоменное взаимодействие.
    \item Тесты решения должны находиться в \texttt{src\_solution/tests/**}.
    \item Сквозной e2e-сценарий должен соответствовать \texttt{docs/operational\_scenario\_v1.md}.
  \end{itemize}
\end{frame}

\begin{frame}{Практический маршрут на 2 дня}
  \begin{itemize}
    \item \textbf{День 1:} структура решения, домены, монитор безопасности, unit/module тесты.
    \item \textbf{День 2:} integration/e2e, сертификация, отчёт, финальный прогон оценки.
    \item Рабочий чек-лист и типичные ошибки смотрите в \texttt{docs/quickstart\_2days.md}.
  \end{itemize}
\end{frame}

\begin{frame}{Обязательные команды самопроверки}
  \begin{itemize}
    \item \texttt{make install}
    \item \texttt{make tests-all}
    \item \texttt{make evaluate-score}
    \item \texttt{make certify-abu-solution}
  \end{itemize}
  \vspace{0.6em}
  \textbf{Платформа проверки:} Linux (или совместимая среда: WSL2/Codespaces/контейнер).
\end{frame}

\begin{frame}{Критерии оценки: как устроена шкала}
  \begin{itemize}
    \item \textbf{25 критериев} (C01--C25), каждый 0--3 балла; номинальная сумма \textbf{75}.
    \item Итог оценки: только \textbf{сумма по критериям (raw)}.
    \item Автоматически: C01--C19, C23--C25; экспертно жюри: C20--C22.
    \item Детали уровней 0--3 --- в \texttt{docs/criteria\_rubric.md}.
  \end{itemize}
\end{frame}

\begin{frame}{Где чаще всего теряются баллы}
  \begin{itemize}
    \item C02: тесты решения размещены вне \texttt{src\_solution/tests/**}.
    \item C07: есть скрипты, но нет успешной сертификации решения.
    \item C08: e2e-тест не покрывает сквозной сценарий из \texttt{docs/operational\_scenario\_v1.md}.
    \item C23--C25: раздутые домены ДВБ, избыточные IPC-интерфейсы, слабые тесты монитора.
  \end{itemize}
\end{frame}

\begin{frame}{Примеры из \texttt{references}}
  \begin{itemize}
    \item \texttt{references/isolation} --- процессная изоляция и базовый IPC.
    \item \texttt{references/secure\_ipc} --- политики маршрутизации и подход ``запрещено по умолчанию''.
    \item \texttt{references/traffic\_light\_demo} --- наглядный учебный сценарий с монитором безопасности.
    \item Краткое руководство по ``Светофору'': \texttt{docs/traffic\_light\_demo.md}.
  \end{itemize}
\end{frame}

\begin{frame}{Финальный чек-лист перед сдачей}
  \begin{itemize}
    \item Пройдены \texttt{make tests-all} и \texttt{make evaluate-score}.
    \item Пройдена \texttt{make certify-abu-solution} с успешным ответом Регулятора.
    \item Отчёт заполнен в \texttt{src\_solution/docs/solution.md}.
    \item Подготовлена сдача по правилам из \texttt{docs/contest\_task.md} и \texttt{docs/contest\_regulations.md}.
  \end{itemize}
\end{frame}

\begin{frame}{Куда смотреть, если застряли}
  \begin{itemize}
    \item Сначала: \texttt{docs/quickstart\_2days.md}.
    \item Затем: \texttt{docs/faq\_participant.md}.
    \item Для формальных деталей: \texttt{docs/criteria\_rubric.md}, \texttt{docs/contest\_regulations.md}.
    \item Для архитектуры и примеров: \texttt{docs/architecture.md}, \texttt{docs/traffic\_light\_demo.md}.
  \end{itemize}
\end{frame}

\begin{frame}{Итог}
  \begin{itemize}
    \item Читайте только \textbf{3 обязательных документа}, остальное --- по необходимости.
    \item Держите фокус на проверяемом результате: тесты, сертификация, отчёт.
    \item Не усложняйте архитектуру сверх необходимого для критериев C01--C25.
  \end{itemize}
  \vspace{0.8em}
  \centering\large Удачи на соревновании!
\end{frame}

\end{document}
